\section{Introduction}
\label{sec:intro}

Tampering with medical images is a demonstrated threat.
In the CT-GAN study, lung nodules injected into or removed from CT scans fooled radiologists in 99\% and 94\% of cases~\cite{mirsky2019ctgan}. % src: docs/council/context_dossier.md (G, CT-GAN row)
Text-guided editors now make such edits easy across modalities~\cite{medbanana2025,chen2026medforge}, and detecting them has become an active benchmark area~\cite{zingarini2024m3dsynth,grabovski2025btd,li2025medforensics,chen2026medforge}.

These benchmarks score a detector by how well it separates untouched images from edited ones.
But a generative edit is two interventions, not one.
A lesion is painted inside a mask, and the image (or a crop of it) is encoded and decoded by the generator's autoencoder.
In existing medical corpora the untouched class never passes through that autoencoder, so a detector can score well by recognising the autoencoder and never looking at the lesion.
For natural images this shortcut is known, and it has been turned into better training data: Aligned Datasets~\cite{rajan2025aligned} trains on autoencoder reconstructions of real images, B-Free~\cite{guillaro2025bfree} on self-conditioned inpaintings with the background restored, and INP-X~\cite{nebioglu2026inpx} shows that detection rates fall from over 91\% to about 55\% once the inpainted image's background is restored. % src: docs/council/round2/irp_prior_art.md (Controls and threat model)
No medical forgery benchmark controls for it.

Whether autoencoder processing alone should count as forgery depends on the definition.
Under a provenance definition, where any image a generative model has touched is synthetic, it should.
We adopt CT-GAN's tamper definition instead~\cite{mirsky2019ctgan}: a forgery is content changed between the scanner and the archive.
Under that definition, an image that passed an autoencoder unedited is authentic, and a provenance detector cannot tell a malicious edit from benign neural processing.

We find that medical forgery detectors, including MedForge's own, respond to autoencoder processing rather than to the edit.
MedForge-Reasoner~\cite{chen2026medforge}, an open-weight detector trained on MedForge-90K, flags all 48 MedForge forgeries we sampled % src: paper/data/medforge_reasoner_hardlabel.csv
and 0.5\% of 186 untouched MedForge test scans. % src: paper/data/medforge_reasoner_hardlabel.csv
After those scans pass through the Stable Diffusion inpainting autoencoder with no edit, it flags 95.8\% of brain MRI and 88.9\% of fundus scans (\cref{fig:medforge}). % src: docs/council/round3/experiments/local/r3_box_audit.csv (N_sd, real, fake_rate)
It is wrong under either definition: every one of those verdicts comes with a lesion-sized box, comparable in size to the boxes it draws on real MedForge forgeries, whereas an image that is synthetic as a whole would call for a whole-image box (\cref{sec:confound}). % src: docs/council/round3/experiments/local/R3_STEP1_REPORT.md (Sec. 1: every null-twin fake call has a box)
The same holds for conventionally trained detectors on our own data.
Mesorch~\cite{zhu2025mesorch} fine-tuned to separate sources from forgeries scores AUC 0.963 on that task % src: work/e0/out_local/tables/table_main.csv
but 0.536 when the negatives are the forgeries' own null twins, % src: work/e0/out_local/tables/table_main.csv
images that passed the same autoencoder without being edited.
By our fingerprint-share metric (\cref{sec:protocol}), 93\% of its margin over chance is reproduced by an image nobody edited. % src: paper/data/fingerprint_share_test.csv (derived); work/e0/out_local/tables/E1_REPORT.md

To make this measurable, we release \dataset.
To our knowledge, it is the first publicly released multi-modality medical local-edit forgery benchmark that pairs every generative edit with a per-pixel (non-box) mask, a background bit-identical to the source, and an autoencoder-matched null twin that passes the same generator without an edit.
Each \emph{lineage} is one source image with one mask and one random seed, from which every other image is derived: the source itself, a \emph{null twin} per generator (the source passed through that generator's autoencoder inside the mask, with nothing edited), the lesion injection or removal, and two copy-move controls.
\dataset holds 3,956 lineages across CT, endoscopy, histopathology, MRI, ultrasound and X-ray, % src: docs/council/context_dossier.md (A)
edited by Stable Diffusion~1.5 inpainting and big-LaMa, with SDXL inpainting held out as a third generator.

AUC alone cannot rank detectors on \dataset, so the protocol fixes one operating point on calibration and reads it against every kind of unedited image.
Measured this way, single-image detection falls from 0.963 on the conventional task % src: work/e0/out_local/tables/table_main.csv
to at most 0.768 against the twin, and stays near chance on CT. % src: work/e0/out_local/tables/e9_final_comparison.csv
Comparing with the null twin as a reference makes AUC easy, but not the operating point.
A zero-parameter pixel subtraction reaches 0.921, % src: work/e0/out_local/e0x/e0a/e0a_deltas.csv
or 0.976 when standardised per modality, % src: docs/council/round3/experiments/local/r3_subtraction_standardised.csv (standardised, ALL)
which is within 0.008 of a learned FC-Siam-diff baseline~\cite{daudt2018fcsiam}; % src: docs/council/round3/experiments/local/r3_subtraction_gap_to_p1n.csv (standardised, ALL)
yet at 95\% recall subtraction flags 47\% of test null twins, % src: docs/council/round3/experiments/local/r3_operating_points.csv (sub_raw, fpr_test_SD_null)
and the learned baseline flags 78\% of denoised sources. % src: docs/council/round3/experiments/local/r3_operating_points.csv (p1n, fpr_nlmeans)
No detector we test is acceptable on every negative type.
Our contributions are:
\begin{itemize}[leftmargin=*,itemsep=1pt,topsep=2pt]
\item \textbf{Dataset.} \dataset, a six-modality local-edit benchmark with per-pixel lesion-shaped masks, bit-identical backgrounds and an autoencoder-matched null twin for every generative edit, over three generators with one held out (\cref{sec:dataset}).
\item \textbf{Protocol.} Four controls that make a forgery-detection number interpretable: an autoencoder-matched unedited twin for every edit; one calibration operating point scored against every negative type, per modality, with same-generator drift beside cross-generator transfer; null-penalised localisation F1; and a mask-shape baseline beside every removal number. These come with frozen leak-free splits and calibration-only selection (\cref{sec:protocol}).
\item \textbf{Findings.} The autoencoder response in three settings, including MedForge's own detector; per-modality operating points for single-image and reference-based detectors, which show that AUC gains do not carry over to false-alarm rates; and an oracle-mask diagnostic separating finding an edit from judging it (\cref{sec:results}).
\end{itemize}
